\section{Билет 67. Термодинамическое определение энтропии. Вычисление приращения энтропии для изопроцессов идеального газа: изотермического и изобарического.}

\begin{definition}
\textbf{Энтропия} $S$ --- функция состояния термодинамической системы, дифференциал которой для обратимого процесса определяется как:
\begin{equation}
    dS = \frac{\delta Q_{\text{обр}}}{T}. \label{eq:entropy_def_67}
\end{equation}
Поскольку $\oint \frac{\delta Q_{\text{обр}}}{T} = 0$, интеграл $\int_1^2 \frac{\delta Q_{\text{обр}}}{T}$ не зависит от пути перехода, а определяется только начальным и конечным состояниями. Это доказывает, что $S$ является функцией состояния.
\end{definition}

\begin{theorem}[Приращение энтропии в изотермическом процессе]
В изотермическом процессе температура постоянна: $T = \text{const}$. Для идеального газа внутренняя энергия зависит только от температуры, поэтому $\Delta U = 0$. Из первого начала термодинамики:
\begin{equation}
    \delta Q = dU + p\,dV = p\,dV.
\end{equation}
Используя уравнение состояния $pV = \nu RT$, получаем:
\begin{equation}
    \delta Q = \frac{\nu RT}{V} dV.
\end{equation}
Подставляя в определение энтропии \eqref{eq:entropy_def_67} и интегрируя от состояния 1 к состоянию 2:
\begin{equation}
    \Delta S_{\text{изот}} = \int_1^2 \frac{\delta Q}{T} = \nu R \int_{V_1}^{V_2} \frac{dV}{V} = \nu R \ln\frac{V_2}{V_1}. \label{eq:delta_S_iso_67}
\end{equation}
Используя закон Бойля--Мариотта $p_1 V_1 = p_2 V_2$, энтропию можно выразить через давления:
\begin{equation}
    \Delta S_{\text{изот}} = \nu R \ln\frac{p_1}{p_2}.
\end{equation}
\end{theorem}

\begin{theorem}[Приращение энтропии в изобарическом процессе]
В изобарическом процессе давление постоянно: $p = \text{const}$. Элементарное количество теплоты:
\begin{equation}
    \delta Q_p = \nu C_p \, dT,
\end{equation}
где $C_p$ --- молярная теплоёмкость при постоянном давлении.
Подставляя в определение энтропии:
\begin{equation}
    dS = \frac{\delta Q_p}{T} = \nu C_p \frac{dT}{T}.
\end{equation}
Интегрируя от $T_1$ до $T_2$:
\begin{equation}
    \Delta S_{\text{изоб}} = \nu C_p \int_{T_1}^{T_2} \frac{dT}{T} = \nu C_p \ln\frac{T_2}{T_1}. \label{eq:delta_S_isob_67}
\end{equation}
Используя закон Гей-Люссака $V/T = \text{const}$, энтропию можно выразить через объёмы:
\begin{equation}
    \Delta S_{\text{изоб}} = \nu C_p \ln\frac{V_2}{V_1}.
\end{equation}
\end{theorem}

\begin{remark}
\textbf{Физическая интерпретация:}
\begin{itemize}
    \item В изотермическом расширении ($V_2 > V_1$) энтропия возрастает ($\Delta S > 0$), так как газ получает теплоту от нагревателя и переходит в состояние с большей неупорядоченностью.
    \item В изобарическом нагревании ($T_2 > T_1$) энтропия также возрастает, поскольку растёт как температура, так и объём газа.
    \item При обратных процессах (сжатие, охлаждение) энтропия системы уменьшается, но энтропия ``система + окружающая среда'' всегда возрастает в соответствии со вторым началом термодинамики.
\end{itemize}
\end{remark}
